\chapter{Literature Review}

\section{General}
This chapter surveys existing work and commonly proposed student projects
in agricultural machine learning, identifies which problem framings are
already saturated, and positions the present project relative to that
landscape.

\section{Commonly Proposed Approaches (and their limitations)}
A survey of recent student-project repositories, dataset platforms, and
government-hackathon problem statements shows that a small number of
problem framings recur very frequently in the "agricultural ML" space:
crop yield prediction and crop-disease image classification (typically
using the PlantVillage or similar Kaggle datasets), and generic
price-prediction notebooks that report an accuracy or error metric
without any accompanying usable interface. These approaches are
well-documented and technically approachable, but by the same token are
heavily repeated -- reviewers and evaluators encounter near-identical
submissions regularly, which reduces their value as a demonstration of
original problem-solving.

Two structural weaknesses recur across this literature:
\begin{itemize}
    \item \textbf{Granularity mismatch.} Many price/yield-prediction
    projects operate at a state or district level, which can mask
    meaningful local variation -- the same crop can differ significantly
    in price between two mandis in the same state on the same day. A
    model trained at too coarse a granularity cannot support a genuinely
    local, actionable recommendation.
    \item \textbf{Missing "last mile".} Most academic implementations stop
    at a trained model and a reported accuracy figure, without an
    interface that would let a non-technical end user (a farmer, in this
    domain) actually act on the model's output.
\end{itemize}

\section{Data Sources for Indian Agricultural Prices}
The Government of India's \textbf{Agmarknet} portal (under data.gov.in)
publishes daily wholesale/modal commodity prices from regulated mandis
across the country, and is the primary live data source used by
government and academic work in this space alike. For historical depth,
this project uses the Kaggle-hosted \textbf{"Daily Market Prices of
Commodity India"} dataset, compiled by the same maintainer as the live
API dataset from the same underlying official Agmarknet records. Using
one consistent data lineage across historical backfill and live
collection avoids the schema and coverage inconsistencies that can arise
from stitching together independently-compiled sources.

\section{Forecasting Methodology}
Time-series forecasting of agricultural commodity prices has been
approached in the literature using classical statistical methods (ARIMA
and its variants), and more recently using additive decomposition models
such as \textbf{Facebook Prophet}, which explicitly models trend and
seasonal components and is well-suited to the kind of small,
per-market datasets this project works with, without requiring extensive
hyperparameter tuning. This project adopts Prophet as its primary
forecasting model, benchmarked against a naive/moving-average baseline
(see Chapter~3).

\section{Positioning of the Present Work}
Relative to the reviewed literature, this project is distinguished by:
(a) forecasting at the individual (commodity, mandi) level rather than a
coarser aggregate, preserving the local price variation that is the
actual decision-relevant signal for a farmer; (b) combining a live,
continuously updating data pipeline with a one-time historical backfill,
rather than a single static training run; and (c) delivering the model's
output through a complete, usable web interface rather than stopping at
a notebook-level result.
